\section{Sovereignty and Security}
\label{sec:sovereignty}

The problem statement is unusually direct about this requirement:

\begin{quote}\itshape
The system should also show, through logs or a visible network monitor, that no
external calls are made at any point. That's the actual proof of the sovereign
claim, not just a statement of it.
\end{quote}

Accordingly, sovereignty is implemented in four independent layers, so that no
single failure silently removes the guarantee.

\subsection{Layer 1 --- kernel network isolation}

Generated code executes inside an unprivileged user and network namespace
created with \code{unshare -rn}. The child process receives a private network
stack containing only a down loopback interface. Supplementary controls: RLIMIT
caps on CPU, address space, process count and file size; a scrubbed environment
with credential-shaped variables removed; a scratch working directory; and a
wall-clock timeout.

\begin{lstlisting}
>>> run_python("import socket; socket.create_connection(('1.1.1.1', 53))")
backend = unshare-netns
stdout  = blocked: OSError [Errno 101] Network is unreachable
\end{lstlisting}

\noindent This is a kernel guarantee. The code is not prevented from calling
out by policy or by a library wrapper; there is no route for the packet to
take. Where Docker is available the sandbox upgrades automatically to
\code{-{}-network none} with a read-only root filesystem.

\subsection{Layer 2 --- host egress denial}

At the venue, \code{scripts/04\_lockdown.sh} installs an nftables table with a
default \code{drop} policy on the output hook, permitting only loopback and
established connections. Ollama, bound to \code{127.0.0.1}, continues to
function; everything else stops.

\begin{lstlisting}
table inet sovereign {
  chain output {
    type filter hook output priority 0; policy drop;
    oif "lo" accept
    ip daddr 127.0.0.0/8 accept
    ct state established,related accept
    counter comment "dropped-egress"
  }
}
\end{lstlisting}

\subsection{Layer 3 --- live monitoring and the canary}

A background sentinel samples connection state and classifies every peer as
local or external, exposing counters through \code{GET /sovereignty} and a
persistent panel in the interface. A deliberate breach attempt is available at
\code{POST /sovereignty/canary}.

\begin{lstlisting}
{"blocked": true, "verdict": "BLOCKED",
 "detail": "TimeoutError: timed out", "ms": 3003}
\end{lstlisting}

\noindent The canary is the most valuable twenty seconds of any demonstration:
it converts an assertion into an observation the audience can watch happen. The
monitor reports honestly in both directions --- on a development machine with
ordinary internet access it correctly reports that the call succeeded.

\subsection{Layer 4 --- tamper-evident audit}

Every prompt, model selection, tool invocation and file write is appended to a
hash chain. Each record embeds the SHA-256 digest of its predecessor, so any
modification or deletion breaks verification and the first corrupted index is
reported.

\begin{lstlisting}
record_n = { ts, event, payload, prev: hash(record_{n-1}) }
hash_n   = sha256(canonical_json(record_n))
\end{lstlisting}

\noindent A regression test alters one field in the middle of a chain and
asserts that verification fails at exactly that index. This makes the audit
log suitable as compliance evidence rather than merely as a debugging aid.

\subsection{Additional controls}

\begin{center}
\renewcommand{\arraystretch}{1.2}
\begin{tabularx}{\linewidth}{@{}l X@{}}
\toprule
\textbf{Control} & \textbf{Implementation} \\
\midrule
Loopback assertion & The LLM client refuses to construct against a non-local host \\
No listening ports & Graph and vector stores are embedded libraries \\
Filesystem jail & All file tools resolve paths and reject workspace escape \\
Offline model cache & \code{HF\_HUB\_OFFLINE}, \code{TRANSFORMERS\_OFFLINE} set \\
Typed tool arguments & Coerced and validated before invocation \\
No error leakage & Failed steps cannot be substituted into deliverables \\
\bottomrule
\end{tabularx}
\end{center}

\subsection{Threat model}

\begin{center}
\renewcommand{\arraystretch}{1.2}
\begin{tabularx}{\linewidth}{@{}X l X@{}}
\toprule
\textbf{Threat} & \textbf{Status} & \textbf{Mitigation} \\
\midrule
Model exfiltrates data over the network & Mitigated & no route exists from the host \\
Generated code opens a socket & Mitigated & network namespace \\
Generated code reads outside the workspace & Mitigated & path jail; sandbox scratch dir \\
A dependency phones home on import & Mitigated & host egress denial; offline flags \\
Audit log altered after the fact & Detected & hash chain verification \\
Prompt injection from an ingested document & \textbf{Partial} & tools are typed and enumerated, but a malicious document could still steer the planner \\
Malicious insider with shell access & \textbf{Out of scope} & requires OS-level controls \\
\bottomrule
\end{tabularx}
\end{center}
